\section{Worked problems: pricing the two phases}

\subsection{A 4,096-token prompt on the laptop}
\textbf{Problem.} Predict the time to first token of a 4,096-token prompt for
Llama~3.2~3B on the M1's GPU from equation~\eqref{eq:prefill-time} and the
measured prefill plateau. Say which term the prediction neglects, estimate it,
and compare both with the measurements of 18.2~s (flash attention off) and
15.7~s (on).

\textbf{Solution.} The plateau of Figure~\ref{fig:ch04-prefill-curve} is about
295 prompt tokens per second, so equation~\eqref{eq:prefill-time}, which charges
every token the same $2P$ FLOPs, predicts $4{,}096 / 295 = 13.9$~s. It neglects
causal attention, whose cost per token grows with the prompt: about
$2S^{2}H_qd_hL = 2 \times 4{,}096^{2} \times 3{,}072 \times 28 = 2.9$~TFLOP,
against $2 \times 2.82 \times 10^{9} \times 4{,}096 = 23.1$~TFLOP for the layers'
matrices, 12.5\% more work. If attention ran as efficiently as the matrices, the
prompt would take $13.9 \times 1.125 = 15.6$~s --- the flash-attention
measurement to within 1\%. The default path took 18.2~s: its attention ran at
about 40\% of the matrices' rate. The prediction is wrong in the optimistic
direction, and by more the longer the prompt, since attention's share grows
linearly with $S$; by about 47,000 tokens for a model of Qwen3-8B's shape it
would be half the work.

\subsection{The same intensity, different costs}
\textbf{Problem.} Show that prefill of $S$ tokens and batched decode of $B$
sequences have the same arithmetic intensity in the linear layers when $S = B$.
List what nevertheless makes them cost different amounts, with numbers for
Llama~3.2~3B at $S = B = 32$ and 4,096 tokens of context per decoding sequence.

\textbf{Solution.} In both cases every weight matrix is multiplied by a matrix
of 32 activation rows: $2 \times 32$ FLOPs per weight read, $2 \times 32 / b_w$
per byte, identical. The differences are elsewhere. (1) \emph{Attention bytes.}
The 32 prompt tokens attend to each other, reading at most 32 tokens of keys
and values; each of the 32 decoding sequences reads its own 4,096 tokens, 470~MB
of cache per sequence and 15~GB per step --- seven and a half times the weights, all of it
at 3~FLOP per byte. (2) \emph{Attention FLOPs.} The prompt's causal attention
is about $2 \times 32^2 \times 3{,}072 \times 28 = 0.18$~GFLOP; the decode
batch's is $32 \times 12{,}288 \times 4{,}096 \times 28 = 45$~GFLOP. (3)
\emph{Output projection.} Prefill needs logits for one position, decode for
32: $32 \times 788$~MFLOP $= 25$~GFLOP against 0.8. (4) \emph{Cache writes.}
Both write 32 tokens' keys and values, 3.7~MB. (5) \emph{Latency.} The prefill
step delivers one token, the first of one request; the decode step delivers one
token to each of 32 requests. The linear layers are the same product; the step
around them is not.

\subsection{Choosing a chunk size in tiles}
\textbf{Problem.} On the M1's GPU the tiled matrix kernel costs about 111~ms per
32-column tile at large sizes, and the measured step times are 122~ms for up to
32 rows, 233~ms for 64 and 325~ms for 66. Two sequences are decoding when a
512-token prompt arrives. For prompt chunks of 30, 62, 64 and 126 tokens,
each processed with the two decoding rows, compute the step time (the worst gap
the decoders see) and the time from the start of the prefill to the prompt's
first token. Compare with processing the prompt in one step.

\textbf{Solution.} One step carries $512 + 2 = 514$ rows, 17 tiles, about
$17 \times 111 = 1.89$~s. Chunks of 62 give steps of 64 rows, two tiles,
0.233~s; eight such steps cover 496 tokens and a ninth carries the last 16
(18 rows, one tile, 0.122~s): first token after $8 \times 0.233 + 0.122 =
1.98$~s. Chunks of 64 give 66 rows, three tiles, 0.325~s per step: $8 \times
0.325 = 2.60$~s. Chunks of 126 give 128 rows, four tiles, 0.432~s per step; four
steps cover 504 tokens and a fifth carries 8 plus the two decode rows (10 rows,
one tile): $4 \times 0.432 + 0.122 = 1.85$~s. Chunks of 30 give 32 rows, one
tile, 0.122~s per step; seventeen steps cover 510 tokens and an eighteenth
carries the last 2 plus two decode rows (4 rows, the small-batch kernel,
71~ms): $17 \times 0.122 + 0.071 = 2.14$~s (all derived from the measured step
times). Summarized as (worst gap, first token): one step (1.89, 1.89); 126
(0.43, 1.85); 62 (0.23, 1.98); 64 (0.32, 2.60); 30 (0.12, 2.14). The 64-token
chunk is dominated by the 62-token chunk on both axes. The general rule: pick
the chunk so that chunk plus decode rows is a multiple of the kernel's tile,
which on this machine means $32k - B_{\text{decode}}$; an engine that knows its
kernel's tile can choose its token budget per step to fit
(\Chref{chunked-prefill}).
